%%%%%%%%%%%%%%%%%%%%%%%%%%%%%%%%%%%%%%%%%%%%%%%%%%%%%%%%%%%%%%%%%%%%%%%%%%%
% Englische Fassung des Uebungsblatts zu Lehreinheit 2 (Teil A und B).
% Vorlage: handouts/LE02.tex. Nur in handouts.tex mit \def\wuzsprache{en}.
%%%%%%%%%%%%%%%%%%%%%%%%%%%%%%%%%%%%%%%%%%%%%%%%%%%%%%%%%%%%%%%%%%%%%%%%%%%
\begin{uebungsblatt}{Dissecting Questions}

\aufgabe{Putting questions to the test}[in class, pairs, 10 min.]
The following questions are constructed. Choose two of them. Note what one
must already believe for the question to make sense (premise), and write a
version that also allows an unexpected answer.

\begin{enumerate}
  \item Why do students learn less in lectures with a laptop than with pen
    and paper?
  \item How strongly does the full moon influence the number of emergency
    admissions?
  \item Which strategies do successful founders use to cope with setbacks?
  \item At what age do children lose their creativity?
  \item What percentage of their brain do people use in everyday life?
  \item Why are women the better leaders?
  \item How can the university prevent ChatGPT from destroying independent
    thinking?
  \item How many hours of sleep does a person need to perform well?
\end{enumerate}

\begin{center}
\renewcommand{\arraystretch}{2.2}
\begin{tabular}{|p{0.06\textwidth}|p{0.38\textwidth}|p{0.46\textwidth}|}
\hline
\textbf{No.} & \textbf{Premise(s)} & \textbf{More open version} \\
\hline
 & & \\ \hline
 & & \\ \hline
\end{tabular}
\end{center}

\aufgabe{Dissecting your own question}[in class, pairs, 20 min.]
Write down the question you brought along. Your partner asks you the
questions from the table; you note the answers. If you have no question of
your own, choose a substitute topic: cafeteria food, commuting to
university, learning with flashcards, part-time jobs during studies.

\smallskip
\noindent\textbf{My question (draft):}
\schreiblinien{2}

\begin{center}
\renewcommand{\arraystretch}{2.0}
\begin{tabular}{|p{0.36\textwidth}|p{0.56\textwidth}|}
\hline
What is presupposed? & \\ \hline
Which words are vague? & \\ \hline
What would I measure (indicator)? & \\ \hline
With whom, where, when? & \\ \hline
Which answer would surprise me? & \\ \hline
\end{tabular}
\end{center}

\aufgabe{SMART check}[in class, pairs]
Rephrase the question and tick the criteria it meets: \emph{specific} (one
object, one population, one context), \emph{measurable} (clear which
quantity is collected), \emph{achievable} (access to data, people,
material), \emph{relevant} (the answer changes knowledge or practice),
\emph{time-bound} (feasible in the available time).
\schreiblinien{2}
\noindent$\square$ specific \quad $\square$ measurable \quad
$\square$ achievable \quad $\square$ relevant \quad $\square$ time-bound
\end{uebungsblatt}

\begin{uebungsblattb}{4}{Recording the Research Question}

\aufgabe{Research question in writing}[homework, due by Unit~3]
Formulate a scientific research question on your semester topic. Explain
in three sentences why this question can be answered and is not trivial.
Also complete the test sentence: \enquote{I am studying \dots, because I
want to find out \dots, in order to help others understand \dots}
\schreiblinien{5}

\aufgabe{Diary of wonder}[homework]
Phrase this week's moment of wonder as a question and note below it the
premise it contains.

\begin{lehrhinweis}[Expected answers]
Task~1, possible premises: (1) that laptop users learn less (direction
presupposed) and that \enquote{learning} is clearly measured; (2) that
there is an influence; (3) that \enquote{successful} is defined and that
success is linked to the strategies (survivorship bias: those who failed
are missing); (4) that children lose creativity and that it can be
measured as one quantity; (5) presupposes a fixed percentage, the
well-known ten-percent myth; (6) direction and group comparison
presupposed, \enquote{better} undefined; (7) that thinking is being
destroyed and that the university can prevent it; (8) that there is one
number for all people and that \enquote{perform well} is clear. More open
versions ask \emph{whether} and \emph{how strongly} rather than
\emph{why}, and name population and measure. Tasks~2/3: in the end the
question must contain population, measure and time frame; a why-question
may remain if its premise is checked first. Task~4: the justification
names a data source and shows that the answer is not obvious.
\end{lehrhinweis}
\end{uebungsblattb}
